\section{Auswertung des Jahresabschlusses 2024}
\label{sec:jahresabschluss}

Grundlage dieses Abschnitts ist der Konzernabschluss der Aumann~AG f\"ur das
Gesch\"aftsjahr 2024. Alle Betr\"age sind in Millionen Euro angegeben, die
Ver\"anderungsraten aus den exakten ver\"offentlichten Werten berechnet.

\subsection{Umsatz und Ver\"anderung zum Vorjahr}
\label{subsec:umsatz}

Der Konzernumsatz belief sich im Gesch\"aftsjahr 2024 auf
\MioEUR{\UmsatzZFIV} nach \MioEUR{\UmsatzZFIII} im
Vorjahr.\quelleGB{2024}{2}\quelleMerken{p6-gb2024-2} Die relative Ver\"anderung ergibt sich als
\begin{equation}
\label{eq:umsatzveraenderung}
\Delta_{\text{Umsatz}}
= \frac{\text{Umsatz}_{2024} - \text{Umsatz}_{2023}}{\text{Umsatz}_{2023}}
\end{equation}
und damit
\[
\Delta_{\text{Umsatz}}
= \frac{\MioEURexakt{\UmsatzZFIV} - \MioEURexakt{\UmsatzZFIII}}
       {\MioEURexakt{\UmsatzZFIII}}
= \Pct{\UmsatzDeltaZFIV}.
\]
Der Umsatz erreichte damit den h\"ochsten Wert der
Unternehmensgeschichte.\quelleGB{2024}{13}

\subsection{Auftragseingang und Ver\"anderung zum Vorjahr}
\label{subsec:auftragseingang}

Der Auftragseingang lag 2024 bei \MioEUR{\AuftragseingangZFIV} gegen\"uber
\MioEUR{\AuftragseingangZFIII} im Vorjahr.\quelleErneut{p6-gb2024-2} Analog zu
Gleichung~\eqref{eq:umsatzveraenderung} ergibt sich
\[
\Delta_{\text{Auftragseingang}}
= \frac{\MioEURexakt{\AuftragseingangZFIV} - \MioEURexakt{\AuftragseingangZFIII}}
       {\MioEURexakt{\AuftragseingangZFIII}}
= \Pct{\AuftragseingangDeltaZFIV}.
\]
Der Auftragseingang entwickelte sich damit gegenl\"aufig zum Umsatz. Als
Ursache nennt der Gesch\"aftsbericht die zur\"uckhaltende Investitionst\"atigkeit
der Automobilhersteller.\quelleGB{2024}{5} Der Auftragsbestand ging im selben
Zeitraum ebenfalls zur\"uck.\quelleErneut{p6-gb2024-2}

\subsection{EBITDA-Marge in Prozent vom Umsatz}
\label{subsec:ebitda-marge}

Die EBITDA-Marge setzt das EBITDA ins Verh\"altnis zum Umsatz:
\begin{equation}
\label{eq:ebitda-marge}
\text{EBITDA-Marge} = \frac{\text{EBITDA}}{\text{Umsatz}}
\end{equation}

Die Aumann~AG weist zwei Gr\"o{\ss}en nebeneinander aus: das berichtete EBITDA
und das um Sondereffekte bereinigte (\emph{adjusted}) EBITDA.\quelleErneut{p6-gb2024-2}
Aus dem berichteten EBITDA von \MioEUR{\EbitdaBerichtetZFIV} folgt
\[
\text{EBITDA-Marge}_{2024}^{\text{berichtet}}
= \frac{\MioEURexakt{\EbitdaBerichtetZFIV}}{\MioEURexakt{\UmsatzZFIV}}
= \Pct{\EbitdaMargeBerichtetZFIV},
\]
aus dem bereinigten EBITDA von \MioEUR{\EbitdaOpZFIV} entsprechend
\[
\text{EBITDA-Marge}_{2024}^{\text{bereinigt}}
= \frac{\MioEURexakt{\EbitdaOpZFIV}}{\MioEURexakt{\UmsatzZFIV}}
= \Pct{\EbitdaMargeZFIV}.
\]
Im Vorjahr betrugen die entsprechenden Werte
\Pct{\EbitdaMargeBerichtetZFIII} beziehungsweise
\Pct{\EbitdaMargeZFIII}.\quelleErneut{p6-gb2024-2}

\subsection{Operatives EBITDA}
\label{subsec:ebitda-operativ}

Als operatives EBITDA wird im Folgenden das von der Gesellschaft
ausgewiesene bereinigte EBITDA verwendet, da es die Ertragskraft des
laufenden Gesch\"afts ohne Sondereffekte abbildet. Es betrug 2024
\MioEUR{\EbitdaOpZFIV} nach \MioEUR{\EbitdaOpZFIII} im Vorjahr, was einer
Steigerung um \Pct{\EbitdaOpDeltaZFIV} entspricht.\quelleGB{2024}{2}\quelleMerken{p7-gb2024-2}

Die Differenz zum berichteten EBITDA von \MioEUR{\EbitdaBerichtetZFIV}
betr\"agt \MioEURexakt{\BereinigungsbetragZFIV} und geht im Wesentlichen auf
die Bereinigung von Personalaufwendungen zur\"uck.\quelleGB{2024}{17}

\begin{table}[htbp]
\centering
\caption{EBITDA-Gr\"o{\ss}en der Aumann~AG 2023 und 2024 in Mio.\,\euro.}
\label{tab:ebitda-vergleich}
\begin{tabular}{@{}lrrr@{}}
\toprule
\textbf{Gr\"o{\ss}e} & \textbf{2023} & \textbf{2024} & \textbf{Marge 2024}\\
\midrule
EBITDA (berichtet)  & \num{\EbitdaBerichtetZFIII} & \num{\EbitdaBerichtetZFIV}
                    & \Pct{\EbitdaMargeBerichtetZFIV}\\
EBITDA (bereinigt)  & \num{\EbitdaOpZFIII}        & \num{\EbitdaOpZFIV}
                    & \Pct{\EbitdaMargeZFIV}\\
\bottomrule
\end{tabular}
\end{table}

\noindent Zur Einordnung: Es ist nicht selbstverst\"andlich, dass das
bereinigte EBITDA \"uber dem berichteten liegt. Im Gesch\"aftsjahr 2025 kehrt
sich das Verh\"altnis um: Dort steht ein berichtetes EBITDA von
\MioEUR{\EbitdaBerichtetZFV} einem bereinigten EBITDA von
\MioEUR{\EbitdaOpZFV} gegen\"uber.\quelleGB{2025}{2}

\subsection{Nettogewinn und Ver\"anderung zum Vorjahr}
\label{subsec:nettogewinn}

Der Konzernjahres\"uberschuss betrug 2024 \MioEUR{\NettogewinnZFIV} nach
\MioEUR{\NettogewinnZFIII} im Vorjahr.\quelleErneut{p7-gb2024-2} Daraus folgt
\[
\Delta_{\text{Nettogewinn}}
= \frac{\MioEURexakt{\NettogewinnZFIV} - \MioEURexakt{\NettogewinnZFIII}}
       {\MioEURexakt{\NettogewinnZFIII}}
= \Pct{\NettogewinnDeltaZFIV}.
\]
Der Nettogewinn hat sich damit gegen\"uber dem Vorjahr mehr als verdoppelt.

\subsection{Dividende und Dividendenrendite}
\label{subsec:dividende}

Vorstand und Aufsichtsrat schlugen f\"ur das Gesch\"aftsjahr 2024 eine
Dividende von \EURje{\DividendeJeAktieZFIV} je dividendenberechtigter Aktie
vor.\quelleGB{2024}{8} Die Dividendenrendite bezogen auf den Aktienkurs zum
Jahresende 2024 ergibt sich als
\begin{equation}
\label{eq:dividendenrendite}
\text{Dividendenrendite}
= \frac{\text{Dividende je Aktie}}{\text{Aktienkurs zum Jahresende}}
\end{equation}
Der XETRA-Schlusskurs des letzten Handelstages des Jahres 2024 betrug
\EURje{\KursSilvesterZFIV}.\quelleGB{2024}{13}\quelleMerken{p7-gb2024-13} Damit gilt
\[
\text{Dividendenrendite}_{2024}
= \frac{\EURje{\DividendeJeAktieZFIV}}{\EURje{\KursSilvesterZFIV}}
= \Pct{\DividendenrenditeZFIV}.
\]

% Nicht nummeriert: Die Aufgabenstellung kennt in Teil 2 die Fragen 2.1
% bis 2.7. Der zusammenfassende Ueberblick ist eine Zugabe und darf die
% Nummerierung nicht verschieben - sonst traegt die Kursentwicklung, die
% Frage 2.7 beantwortet, die Nummer 2.8.
% Ueberschrift, Ankuendigungssatz und Tabelle gehoeren zusammen; ohne
% \needspace stand die Ueberschrift samt einer einzelnen Zeile am Fuss von
% Seite 7, waehrend die Tabelle erst auf Seite 8 folgte. Der Wert ist fuer
% alle drei bemessen (Ueberschrift, Satz, Tabellenkoerper mit Beschriftung),
% nicht nur fuer die Ueberschrift - sonst passt genau diese hinein und die
% Tabelle bleibt trotzdem zurueck.
\needspace{16\baselineskip}
\subsection*{\"Uberblick der Kennzahlen 2023 und 2024}
\addcontentsline{toc}{subsection}{\"Uberblick der Kennzahlen 2023 und 2024}

Tabelle~\ref{tab:kennzahlen-ueberblick} fasst die vier in den
Abschnitten~\ref{subsec:umsatz} bis \ref{subsec:nettogewinn} einzeln
hergeleiteten Gr\"o{\ss}en zusammen.

\begin{table}[!htbp]
\centering
\caption{Ausgew\"ahlte Konzernkennzahlen der Aumann~AG. Betr\"age in
Mio.\,\euro, Ver\"anderungen in Prozent.}
\label{tab:kennzahlen-ueberblick}
\begin{tabular}{@{}lrrr@{}}
\toprule
\textbf{Kennzahl} & \textbf{2023} & \textbf{2024} & \textbf{Ver\"anderung}\\
\midrule
Umsatz                    & \num{\UmsatzZFIII}          & \num{\UmsatzZFIV}
                          & \Pct{\UmsatzDeltaZFIV}\\
Auftragseingang           & \num{\AuftragseingangZFIII} & \num{\AuftragseingangZFIV}
                          & \Pct{\AuftragseingangDeltaZFIV}\\
Operatives EBITDA         & \num{\EbitdaOpZFIII}        & \num{\EbitdaOpZFIV}
                          & \Pct{\EbitdaOpDeltaZFIV}\\
Konzernjahres\"uberschuss & \num{\NettogewinnZFIII}     & \num{\NettogewinnZFIV}
                          & \Pct{\NettogewinnDeltaZFIV}\\
\bottomrule
\end{tabular}
\end{table}

\subsection{Entwicklung des Aktienkurses 2022 bis 2025}
\label{subsec:kursverlauf}

Die Aumann-Aktie (ISIN
\ISIN\quelleWeb{Aumann AG, Investor Relations -- Shares}{quelle:shares}{29.07.2026})
notiert seit M\"arz 2017 im Prime Standard der Frankfurter
Wertpapierb\"orse.\quelleGB{2024}{13}\quelleMerken{p8-gb2024-13}
Abbildung~\ref{fig:kursverlauf} zeichnet den Kurspfad vom letzten Handelstag
2021 bis zum letzten Handelstag 2025 nach. Ausgangspunkt ist der
Schlusskurs 2021; danach folgt f\"ur jedes Gesch\"aftsjahr das Jahreshoch,
das Jahrestief und der Jahresschlusskurs, angeordnet nach dem Tag, an dem
der Wert eingetreten ist. Die zugeh\"origen Kurse stehen in
Tabelle~\ref{tab:kursextrema}.

Zwei Eigenheiten der Darstellung sind zu beachten. Erstens ist die
Reihenfolge innerhalb eines Jahres nicht fest: 2022 und 2024 erreichte die
Aktie zuerst ihr Hoch, 2023 und 2025 dagegen zuerst ihr Tief. Zweitens ist
der letzte Handelstag eines Jahres der 29.\ oder 30.~Dezember, nicht der
31.~Dezember; auch die Gesch\"aftsberichte beziehen sich auf den letzten
Handelstag.\quelleErneut{p8-gb2024-13}

\begin{figure}[htbp]
\centering
\begin{tikzpicture}
\begin{axis}[
  width=0.9\linewidth, height=6cm,
  xlabel={Gesch\"aftsjahr}, ylabel={XETRA-Schlusskurs in \euro},
  xtick={2022,2023,2024,2025,2026},
  xmin=2021.9, xmax=2026.1,
  ymin=0, ymax=22, grid=major, grid style={gray!25},
]
% Erst die Verbindungslinie ueber alle Stuetzpunkte, dann dieselben Punkte
% noch einmal als Marken. Beide lesen dieselbe, nach Datum geordnete Datei;
% die Markenform ergibt sich aus der Spalte "art" (Stil kurspunkte in
% preamble.tex).
\addplot[blue!60!black, thick, mark=none]
  table[col sep=comma, x=x, y=kurs, comment chars={\#}]
  {data/kursverlauf.csv};
\addplot[kurspunkte, forget plot]
  table[col sep=comma, x=x, y=kurs, meta=art, comment chars={\#}]
  {data/kursverlauf.csv};
\end{axis}
\end{tikzpicture}
\caption{Kurspfad der Aumann-Aktie im XETRA-Handel vom letzten Handelstag
2021 bis zum letzten Handelstag 2025. Eingezeichnet sind je Gesch\"aftsjahr
das Jahreshoch~(\textcolor{kurshoch}{$\blacktriangle$}), das
Jahrestief~(\textcolor{kurstief}{$\blacktriangledown$}) und der
Jahresschlusskurs~(\textcolor{kursschluss}{$\bullet$}), geordnet nach dem Tag
ihres Eintretens. Die Verbindungslinie ist eine Hilfslinie zwischen diesen
St\"utzpunkten; sie gibt den Verlauf dazwischen nicht wieder.}
\label{fig:kursverlauf}
\end{figure}

Die Jahresschlusskurse betrugen \EURje{\KursSilvesterZFI} (2021),
\EURje{\KursSilvesterZFII} (2022), \EURje{\KursSilvesterZFIII} (2023),
\EURje{\KursSilvesterZFIV} (2024) und
\EURje{\KursSilvesterZFV} (2025).\quelleGB{2023}{14}\quelleGB{2024}{13}\quelleMerken{p9-gb2024-13}\quelleGB{2025}{13}
Daraus ergeben sich Jahresver\"anderungen von \Pct{\KursDeltaZFII} (2022),
\Pct{\KursDeltaZFIII} (2023), \Pct{\KursDeltaZFIV} (2024) und
\Pct{\KursDeltaZFV} (2025). \"Uber den Gesamtzeitraum liegt der Schlusskurs
2025 damit unterhalb des Schlusskurses 2021, mit einem ausgepr\"agten Hoch
zum Jahresende 2023 und einem Einbruch im Jahr 2024.

Die Jahresschlusskurse allein verdecken die Schwankungsbreite innerhalb der
Jahre. Tabelle~\ref{tab:kursextrema} stellt ihnen die H\"ochst- und
Tiefstkurse gegen\"uber und nennt den Handelstag, an dem sie eintraten. Die
Tabelle ist damit zugleich die genaue Lesart der Abbildung: Diese kann drei
ihrer St\"utzpunkte nicht aufl\"osen, weil das Hoch 2023, der Schlusskurs
2023 und das Hoch 2024 innerhalb von f\"unf Handelstagen um den Jahreswechsel
lagen und auf einer Achse \"uber vier Jahre zu einer einzigen Spitze
zusammenfallen.

\begin{table}[!htbp]
\centering
\caption{H\"ochst-, Tiefst- und Schlusskurse der Aumann-Aktie im
XETRA-Handel, mit dem Handelstag des jeweiligen Extremwerts.}
\label{tab:kursextrema}
\begin{tabular}{@{}lrlrlr@{}}
\toprule
\textbf{Jahr} & \textbf{H\"ochstkurs} & \textbf{am}
              & \textbf{Tiefstkurs}  & \textbf{am}
              & \textbf{Schlusskurs}\\
\midrule
2022 & \EURje{\KursHochZFII}  & \DatumHochZFII  & \EURje{\KursTiefZFII}
     & \DatumTiefZFII  & \EURje{\KursSilvesterZFII}\\
2023 & \EURje{\KursHochZFIII} & \DatumHochZFIII & \EURje{\KursTiefZFIII}
     & \DatumTiefZFIII & \EURje{\KursSilvesterZFIII}\\
2024 & \EURje{\KursHochZFIV}  & \DatumHochZFIV  & \EURje{\KursTiefZFIV}
     & \DatumTiefZFIV  & \EURje{\KursSilvesterZFIV}\\
2025 & \EURje{\KursHochZFV}   & \DatumHochZFV   & \EURje{\KursTiefZFV}
     & \DatumTiefZFV   & \EURje{\KursSilvesterZFV}\\
\bottomrule
\end{tabular}
\end{table}

\noindent Das Jahr 2024 tr\"agt dabei die gr\"o{\ss}te Spanne: dem Hoch von
\EURje{\KursHochZFIV} im Januar steht ein Tief von \EURje{\KursTiefZFIV} im
November gegen\"uber, ein R\"uckgang um
\Pct{\KursSpanneZFIV}.\quelleWeb{Yahoo Finance, Kurshistorie der Aumann-Aktie, Symbol AAG.DE}{quelle:yahoo}{29.07.2026}\quelleMerken{kurs-yahoo}

\medskip\noindent
\emph{Hinweis zur Datengrundlage:} Die Gesch\"aftsberichte ver\"offentlichen
weder H\"ochst- noch Tiefstkurse, und die Quartalsmitteilungen enthalten keine
Kursangaben. Die Extremwerte stammen deshalb aus der Kurshistorie von
Yahoo~Finance.\quelleErneut{kurs-yahoo} Diese Reihe ist gegen die
Prim\"arquelle gepr\"uft: F\"ur 2022 bis 2025 stimmen ihre
Dezember-Schlusskurse auf den Cent mit den XETRA-Jahresschlusskursen der
Gesch\"aftsberichte \"uberein, und die Schnittstelle weist XETRA als
Handelsplatz aus.

Zwei Einschr\"ankungen sind zu nennen. Erstens sind H\"ochst- und Tiefstkurs
\emph{Intraday}-Werte, nicht Extrema der Schlusskurse; sie liegen daher
au{\ss}erhalb der Schlusskursreihe. Zweitens l\"asst sich allein der
Schlusskurs 2021 nicht gegen eine Prim\"arquelle pr\"ufen: Der
Gesch\"aftsbericht 2022 enth\"alt, anders als die Berichte 2023 bis 2025,
\"uberhaupt keine Angabe zum Jahresschlusskurs, und der
Gesch\"aftsbericht 2023 f\"uhrt als Vorjahreswert bereits 2022. Der Wert
f\"ur 2021 st\"utzt sich damit auf dieselbe Reihe, deren \"ubrige vier Jahre
die Gegenprobe bestehen; gepr\"uft ist er selbst aber nicht.
